\documentclass[11pt,a4paper]{article}
\usepackage[utf8]{inputenc}
\usepackage[T1]{fontenc}
\usepackage[spanish,es-minimal,es-noshorthands]{babel}
\usepackage{lmodern}
\usepackage[sfdefault,lining]{FiraSans}
\usepackage[a4paper,top=2.3cm,bottom=2.2cm,left=1.8cm,right=1.8cm,headheight=1.3cm,headsep=0.4cm,footskip=1cm]{geometry}
\usepackage{amsmath,amssymb}
\usepackage{icomma}
\usepackage{graphicx}
\usepackage[table]{xcolor}
\usepackage{tikz}
\usepackage{fontawesome5}
\usepackage{ccicons}
\usepackage{fancyhdr}
\usepackage{lastpage}
\usepackage{tasks}
\usepackage{needspace}
\usepackage[most]{tcolorbox}
\usepackage{array}
\usepackage[hidelinks]{hyperref}
\AtBeginDocument{\hypersetup{pdftitle={Soluciones de la ficha de la unidad 9: Perímetros y áreas. Teorema de Pitágoras. 3º ESO},pdfsubject={Matemáticas 3º ESO},pdfkeywords={teorema de Pitágoras, hipotenusa, catetos, perímetro, área, polígonos regulares, apotema, circunferencia, círculo, sector circular, corona circular}}}

\definecolor{navy}{HTML}{1B2748}
\definecolor{accent}{HTML}{E8603A}
\definecolor{teal}{HTML}{1F9E90}
\definecolor{soft}{HTML}{EEF1F8}
\definecolor{muted}{HTML}{5D6685}

\setlength{\parindent}{0pt}
\setlength{\parskip}{3pt}
\settasks{label=\alph*), label-format=\bfseries\color{navy}, label-width=1.3em, item-indent=1.8em, column-sep=0.6em, after-item-skip=5pt, before-skip=3pt, after-skip=4pt}

% ---------- capçalera i peu ----------
\newcommand{\llicencia}{\href{https://creativecommons.org/licenses/by-nc-sa/4.0/deed.es}{\ccbyncsa}\ \ Licencia Creative Commons Atribución-NoComercial-CompartirIgual 4.0 Internacional (CC BY-NC-SA 4.0)}
\pagestyle{fancy}
\fancyhf{}
\fancyhead[L]{\raisebox{-0.35\height}{\includegraphics[height=0.95cm]{logo-gva.png}}}
\fancyhead[C]{\small\color{muted}Soluciones de la ficha · Unidad 9: Perímetros y áreas. Teorema de Pitágoras · 3\textsuperscript{o} ESO}
\fancyhead[R]{\raisebox{-0.35\height}{\includegraphics[height=0.95cm]{logo-vermellar.png}}}
\fancyfoot[L]{\scriptsize\color{muted}\llicencia}
\fancyfoot[R]{\small\color{muted}\thepage}
\renewcommand{\headrulewidth}{0.4pt}
\renewcommand{\headrule}{\color{muted!40}\hrule height \headrulewidth}
\fancypagestyle{primera}{\fancyhf{}\renewcommand{\headrulewidth}{0pt}
  \fancyfoot[L]{\scriptsize\color{muted}\llicencia}
  \fancyfoot[R]{\small\color{muted}\thepage}}

% ---------- nivells i exercicis ----------
\newcommand{\estrelles}[1]{{\color{accent}\ifcase#1\or\faStar\faStar[regular]\faStar[regular]\or\faStar\faStar\faStar[regular]\or\faStar\faStar\faStar\fi}}
\newcounter{exer}
\newcommand{\ex}[1]{\par\needspace{4\baselineskip}\medskip\stepcounter{exer}\noindent{\color{navy}\bfseries\theexer.}\ \ {\small\estrelles{#1}}\ \ }

\newcommand{\seccio}[2]{\par\needspace{10\baselineskip}\bigskip
  \begin{tcolorbox}[colback=navy,colframe=navy,coltext=white,arc=2mm,boxrule=0pt,left=3mm,right=3mm,top=1.2mm,bottom=1.2mm]
  {\large\bfseries #1\quad #2}
  \end{tcolorbox}\vspace{-2pt}}

\newcolumntype{C}[1]{>{\centering\arraybackslash}p{#1}}
\usepackage{pifont}

% quadre de recordatori (només a la fitxa d'exercicis)
\newcommand{\recorda}[1]{\par\smallskip\begin{tcolorbox}[colback=accent!7,colframe=accent!55,arc=2mm,boxrule=0.5pt,left=2.5mm,right=2.5mm,top=1mm,bottom=1mm]\small{\color{accent}\faLightbulb[regular]}\ \textbf{Recuerda:}\ #1\end{tcolorbox}}
\newcommand{\res}[1]{\boxed{#1}}
\newcommand{\rest}[1]{\fbox{#1}}

\def\fmneed{14}
% =====================================================================
% Capçalera FentMates de les fitxes: un color per a cada curs.
% Els colors són els de la pàgina inicial (LEVELS a assets/js/materials.js);
% el logo ve de _eines/logo/pdf/ (python3 _eines/logo/genera.py).
%
% Ús, al final del preàmbul de la unitat:
%   % =====================================================================
% Capçalera FentMates de les fitxes: un color per a cada curs.
% Els colors són els de la pàgina inicial (LEVELS a assets/js/materials.js);
% el logo ve de _eines/logo/pdf/ (python3 _eines/logo/genera.py).
%
% Ús, al final del preàmbul de la unitat:
%   % =====================================================================
% Capçalera FentMates de les fitxes: un color per a cada curs.
% Els colors són els de la pàgina inicial (LEVELS a assets/js/materials.js);
% el logo ve de _eines/logo/pdf/ (python3 _eines/logo/genera.py).
%
% Ús, al final del preàmbul de la unitat:
%   \input{../comu/fentmates.tex}
%   \fmcurs{1eso}
% i, al començament del document:
%   \capcalera{Matemàtiques}{1r ESO}{Fitxa d'exercicis · Unitat 2}{Divisibilitat}
% Opcional abans de \input: \def\fmneed{14} (espai mínim abans de cada secció).
% =====================================================================
\makeatletter
\def\fm@logodir{../../logo/pdf}
\providecommand{\fmneed}{10}
\def\fm@cursnom{}

% clau del curs, color i color del text damunt del color
\newcommand{\fm@curs}[3]{\@namedef{fm@color@#1}{#2}\@namedef{fm@text@#1}{#3}}
\fm@curs{1eso}{E6B84A}{navy}
\fm@curs{2eso}{3B82F6}{white}
\fm@curs{3eso}{FF7A50}{white}
\fm@curs{4eso-a}{A78BFA}{white}
\fm@curs{4eso-b}{6D28D9}{white}
\fm@curs{1batx-mi}{9B1C1C}{white}
\fm@curs{1batx-ccss}{EF6B6B}{white}
\fm@curs{2batx-mii}{1E6B47}{white}
\fm@curs{2batx-ccss}{2BB3A3}{white}

\newcommand{\fmcurs}[1]{%
  \@ifundefined{fm@color@#1}{\PackageError{fentmates}{Curs desconegut: #1}{}}{}%
  \edef\fm@tmp{\noexpand\definecolor{curs}{HTML}{\@nameuse{fm@color@#1}}%
               \noexpand\colorlet{curstext}{\@nameuse{fm@text@#1}}}\fm@tmp
  \colorlet{cursfosc}{curs!72!black}%
  \def\fm@logo{\fm@logodir/vertical-#1.pdf}%
  \def\fm@marca{\fm@logodir/marca-#1.pdf}}

% ---------- capçalera de la primera pàgina ----------
% #1 matèria (petit), #2 curs, #3 tipus de document i unitat, #4 títol
\newcommand{\capcalera}[4]{%
  \gdef\fm@cursnom{#2}%
  \par\noindent
  \begin{tcolorbox}[colback=curs,colframe=curs,boxrule=0pt,arc=2.2mm,width=3.5cm,height=2.15cm,
      valign=center,left=3mm,right=2mm,top=1mm,bottom=1mm,nobeforeafter,box align=center]
    \color{curstext}{\footnotesize #1}\par\vspace{1pt}{\fontsize{19}{21}\selectfont\bfseries #2\par}%
  \end{tcolorbox}\hspace{5mm}%
  \begin{minipage}[c]{\dimexpr\linewidth-3.5cm-5mm-2.9cm\relax}
    {\color{cursfosc}\small\bfseries\MakeUppercase{#3}}\par\vspace{2pt}
    {\color{navy}\fontsize{20}{23}\selectfont\bfseries #4\par}
  \end{minipage}\hfill
  \begin{minipage}[c]{2.6cm}\raggedleft\includegraphics[width=2.6cm]{\fm@logo}\end{minipage}\par
  \vspace{2pt}}

% ---------- barres de secció del color del curs ----------
\renewcommand{\seccio}[2]{\par\needspace{\fmneed\baselineskip}\bigskip
  \begin{tcolorbox}[colback=curs,colframe=curs,coltext=curstext,arc=2mm,boxrule=0pt,left=3mm,right=3mm,top=1.2mm,bottom=1.2mm]
  {\large\bfseries #1\quad #2}
  \end{tcolorbox}\vspace{-2pt}}

% ---------- estrelles de dificultat del color del curs (a Batxillerat no se'n posen) ----------
\@ifundefined{estrelles}{}{%
  \renewcommand{\estrelles}[1]{{\color{curs}\ifcase#1\or\faStar\faStar[regular]\faStar[regular]\or\faStar\faStar\faStar[regular]\or\faStar\faStar\faStar\fi}}}

% ---------- capçalera i peu de les pàgines ----------
% primera pàgina: logos institucionals al peu
\fancypagestyle{primera}{\fancyhf{}\renewcommand{\headrulewidth}{0pt}
  \fancyfoot[L]{\raisebox{-0.32\height}{\includegraphics[height=0.78cm]{logo-gva.png}}\hspace{5mm}%
    \raisebox{-0.32\height}{\includegraphics[height=0.78cm]{logo-vermellar.png}}}
  \fancyfoot[R]{\scriptsize\color{muted}\href{https://creativecommons.org/licenses/by-nc-sa/4.0/deed.ca}{\ccbyncsa}\ \ CC BY-NC-SA 4.0 · fentmates.org\qquad\small\thepage}}
% la resta: el símbol a l'esquerra i el curs a la dreta, amb el títol al centre com abans
\fancyhead[L]{\raisebox{-0.3\height}{\includegraphics[height=0.62cm]{\fm@marca}}}
\fancyhead[R]{\small\bfseries\color{cursfosc}\fm@cursnom}
\renewcommand{\headrule}{\color{curs!45}\hrule height \headrulewidth}
\makeatother

%   \fmcurs{1eso}
% i, al començament del document:
%   \capcalera{Matemàtiques}{1r ESO}{Fitxa d'exercicis · Unitat 2}{Divisibilitat}
% Opcional abans de \input: \def\fmneed{14} (espai mínim abans de cada secció).
% =====================================================================
\makeatletter
\def\fm@logodir{../../logo/pdf}
\providecommand{\fmneed}{10}
\def\fm@cursnom{}

% clau del curs, color i color del text damunt del color
\newcommand{\fm@curs}[3]{\@namedef{fm@color@#1}{#2}\@namedef{fm@text@#1}{#3}}
\fm@curs{1eso}{E6B84A}{navy}
\fm@curs{2eso}{3B82F6}{white}
\fm@curs{3eso}{FF7A50}{white}
\fm@curs{4eso-a}{A78BFA}{white}
\fm@curs{4eso-b}{6D28D9}{white}
\fm@curs{1batx-mi}{9B1C1C}{white}
\fm@curs{1batx-ccss}{EF6B6B}{white}
\fm@curs{2batx-mii}{1E6B47}{white}
\fm@curs{2batx-ccss}{2BB3A3}{white}

\newcommand{\fmcurs}[1]{%
  \@ifundefined{fm@color@#1}{\PackageError{fentmates}{Curs desconegut: #1}{}}{}%
  \edef\fm@tmp{\noexpand\definecolor{curs}{HTML}{\@nameuse{fm@color@#1}}%
               \noexpand\colorlet{curstext}{\@nameuse{fm@text@#1}}}\fm@tmp
  \colorlet{cursfosc}{curs!72!black}%
  \def\fm@logo{\fm@logodir/vertical-#1.pdf}%
  \def\fm@marca{\fm@logodir/marca-#1.pdf}}

% ---------- capçalera de la primera pàgina ----------
% #1 matèria (petit), #2 curs, #3 tipus de document i unitat, #4 títol
\newcommand{\capcalera}[4]{%
  \gdef\fm@cursnom{#2}%
  \par\noindent
  \begin{tcolorbox}[colback=curs,colframe=curs,boxrule=0pt,arc=2.2mm,width=3.5cm,height=2.15cm,
      valign=center,left=3mm,right=2mm,top=1mm,bottom=1mm,nobeforeafter,box align=center]
    \color{curstext}{\footnotesize #1}\par\vspace{1pt}{\fontsize{19}{21}\selectfont\bfseries #2\par}%
  \end{tcolorbox}\hspace{5mm}%
  \begin{minipage}[c]{\dimexpr\linewidth-3.5cm-5mm-2.9cm\relax}
    {\color{cursfosc}\small\bfseries\MakeUppercase{#3}}\par\vspace{2pt}
    {\color{navy}\fontsize{20}{23}\selectfont\bfseries #4\par}
  \end{minipage}\hfill
  \begin{minipage}[c]{2.6cm}\raggedleft\includegraphics[width=2.6cm]{\fm@logo}\end{minipage}\par
  \vspace{2pt}}

% ---------- barres de secció del color del curs ----------
\renewcommand{\seccio}[2]{\par\needspace{\fmneed\baselineskip}\bigskip
  \begin{tcolorbox}[colback=curs,colframe=curs,coltext=curstext,arc=2mm,boxrule=0pt,left=3mm,right=3mm,top=1.2mm,bottom=1.2mm]
  {\large\bfseries #1\quad #2}
  \end{tcolorbox}\vspace{-2pt}}

% ---------- estrelles de dificultat del color del curs (a Batxillerat no se'n posen) ----------
\@ifundefined{estrelles}{}{%
  \renewcommand{\estrelles}[1]{{\color{curs}\ifcase#1\or\faStar\faStar[regular]\faStar[regular]\or\faStar\faStar\faStar[regular]\or\faStar\faStar\faStar\fi}}}

% ---------- capçalera i peu de les pàgines ----------
% primera pàgina: logos institucionals al peu
\fancypagestyle{primera}{\fancyhf{}\renewcommand{\headrulewidth}{0pt}
  \fancyfoot[L]{\raisebox{-0.32\height}{\includegraphics[height=0.78cm]{logo-gva.png}}\hspace{5mm}%
    \raisebox{-0.32\height}{\includegraphics[height=0.78cm]{logo-vermellar.png}}}
  \fancyfoot[R]{\scriptsize\color{muted}\href{https://creativecommons.org/licenses/by-nc-sa/4.0/deed.ca}{\ccbyncsa}\ \ CC BY-NC-SA 4.0 · fentmates.org\qquad\small\thepage}}
% la resta: el símbol a l'esquerra i el curs a la dreta, amb el títol al centre com abans
\fancyhead[L]{\raisebox{-0.3\height}{\includegraphics[height=0.62cm]{\fm@marca}}}
\fancyhead[R]{\small\bfseries\color{cursfosc}\fm@cursnom}
\renewcommand{\headrule}{\color{curs!45}\hrule height \headrulewidth}
\makeatother

%   \fmcurs{1eso}
% i, al començament del document:
%   \capcalera{Matemàtiques}{1r ESO}{Fitxa d'exercicis · Unitat 2}{Divisibilitat}
% Opcional abans de \input: \def\fmneed{14} (espai mínim abans de cada secció).
% =====================================================================
\makeatletter
\def\fm@logodir{../../logo/pdf}
\providecommand{\fmneed}{10}
\def\fm@cursnom{}

% clau del curs, color i color del text damunt del color
\newcommand{\fm@curs}[3]{\@namedef{fm@color@#1}{#2}\@namedef{fm@text@#1}{#3}}
\fm@curs{1eso}{E6B84A}{navy}
\fm@curs{2eso}{3B82F6}{white}
\fm@curs{3eso}{FF7A50}{white}
\fm@curs{4eso-a}{A78BFA}{white}
\fm@curs{4eso-b}{6D28D9}{white}
\fm@curs{1batx-mi}{9B1C1C}{white}
\fm@curs{1batx-ccss}{EF6B6B}{white}
\fm@curs{2batx-mii}{1E6B47}{white}
\fm@curs{2batx-ccss}{2BB3A3}{white}

\newcommand{\fmcurs}[1]{%
  \@ifundefined{fm@color@#1}{\PackageError{fentmates}{Curs desconegut: #1}{}}{}%
  \edef\fm@tmp{\noexpand\definecolor{curs}{HTML}{\@nameuse{fm@color@#1}}%
               \noexpand\colorlet{curstext}{\@nameuse{fm@text@#1}}}\fm@tmp
  \colorlet{cursfosc}{curs!72!black}%
  \def\fm@logo{\fm@logodir/vertical-#1.pdf}%
  \def\fm@marca{\fm@logodir/marca-#1.pdf}}

% ---------- capçalera de la primera pàgina ----------
% #1 matèria (petit), #2 curs, #3 tipus de document i unitat, #4 títol
\newcommand{\capcalera}[4]{%
  \gdef\fm@cursnom{#2}%
  \par\noindent
  \begin{tcolorbox}[colback=curs,colframe=curs,boxrule=0pt,arc=2.2mm,width=3.5cm,height=2.15cm,
      valign=center,left=3mm,right=2mm,top=1mm,bottom=1mm,nobeforeafter,box align=center]
    \color{curstext}{\footnotesize #1}\par\vspace{1pt}{\fontsize{19}{21}\selectfont\bfseries #2\par}%
  \end{tcolorbox}\hspace{5mm}%
  \begin{minipage}[c]{\dimexpr\linewidth-3.5cm-5mm-2.9cm\relax}
    {\color{cursfosc}\small\bfseries\MakeUppercase{#3}}\par\vspace{2pt}
    {\color{navy}\fontsize{20}{23}\selectfont\bfseries #4\par}
  \end{minipage}\hfill
  \begin{minipage}[c]{2.6cm}\raggedleft\includegraphics[width=2.6cm]{\fm@logo}\end{minipage}\par
  \vspace{2pt}}

% ---------- barres de secció del color del curs ----------
\renewcommand{\seccio}[2]{\par\needspace{\fmneed\baselineskip}\bigskip
  \begin{tcolorbox}[colback=curs,colframe=curs,coltext=curstext,arc=2mm,boxrule=0pt,left=3mm,right=3mm,top=1.2mm,bottom=1.2mm]
  {\large\bfseries #1\quad #2}
  \end{tcolorbox}\vspace{-2pt}}

% ---------- estrelles de dificultat del color del curs (a Batxillerat no se'n posen) ----------
\@ifundefined{estrelles}{}{%
  \renewcommand{\estrelles}[1]{{\color{curs}\ifcase#1\or\faStar\faStar[regular]\faStar[regular]\or\faStar\faStar\faStar[regular]\or\faStar\faStar\faStar\fi}}}

% ---------- capçalera i peu de les pàgines ----------
% primera pàgina: logos institucionals al peu
\fancypagestyle{primera}{\fancyhf{}\renewcommand{\headrulewidth}{0pt}
  \fancyfoot[L]{\raisebox{-0.32\height}{\includegraphics[height=0.78cm]{logo-gva.png}}\hspace{5mm}%
    \raisebox{-0.32\height}{\includegraphics[height=0.78cm]{logo-vermellar.png}}}
  \fancyfoot[R]{\scriptsize\color{muted}\href{https://creativecommons.org/licenses/by-nc-sa/4.0/deed.ca}{\ccbyncsa}\ \ CC BY-NC-SA 4.0 · fentmates.org\qquad\small\thepage}}
% la resta: el símbol a l'esquerra i el curs a la dreta, amb el títol al centre com abans
\fancyhead[L]{\raisebox{-0.3\height}{\includegraphics[height=0.62cm]{\fm@marca}}}
\fancyhead[R]{\small\bfseries\color{cursfosc}\fm@cursnom}
\renewcommand{\headrule}{\color{curs!45}\hrule height \headrulewidth}
\makeatother

\fmcurs{3eso}

\begin{document}
\thispagestyle{primera}

\capcalera{Matemáticas}{3\textsuperscript{o} ESO}{Soluciones de la ficha · Unidad 9}{Perímetros y áreas. Teorema de Pitágoras}
\vspace{8pt}

\begin{tcolorbox}[colback=soft,colframe=soft,arc=2mm,left=3mm,right=3mm,top=1.5mm,bottom=1.5mm]
\small Intenta resolver cada ejercicio antes de mirar su solución. Si no te sale el mismo resultado, repasa los pasos o consulta los apuntes interactivos de la unidad: encontrarás una calculadora de Pitágoras y herramientas que calculan áreas y perímetros paso a paso. Los resultados aproximados están redondeados a las centésimas con la tecla $\pi$ de la calculadora.
\end{tcolorbox}

% =====================================================
\seccio{1}{El teorema de Pitágoras}

\ex{1} \begin{tasks}(3)
  \task $\sqrt{100}=\res{10\text{ cm}}$
  \task $\sqrt{169}=\res{13\text{ cm}}$
  \task $\sqrt{1681}=\res{41\text{ m}}$
  \task $\sqrt{58}\approx\res{7{,}62\text{ cm}}$
  \task $\sqrt{6{,}25}=\res{2{,}5\text{ m}}$
  \task $\sqrt{200}\approx\res{14{,}14\text{ cm}}$
\end{tasks}

\ex{1} \begin{tasks}(3)
  \task $\sqrt{576}=\res{24}$
  \task $\sqrt{225}=\res{15}$
  \task $\sqrt{36}=\res{6}$
  \task $\sqrt{119}\approx\res{10{,}91}$
  \task $\sqrt{36}=\res{6}$
  \task $\sqrt{65}\approx\res{8{,}06}$
\end{tasks}

\ex{1} a) $x=\sqrt{144+256}=\res{20}$ \qquad b) $x=\sqrt{169-25}=\res{12}$ \qquad c) $x=\sqrt{289-225}=\res{8}$

\ex{2} La escalera es la hipotenusa: $h=\sqrt{6{,}5^2-2{,}5^2}=\sqrt{42{,}25-6{,}25}=\sqrt{36}=\res{6\text{ m}}$.

\ex{2} $h=\sqrt{50^2-30^2}=\sqrt{1600}=\res{40\text{ m}}$.

\ex{2} Los dos tramos son perpendiculares: $d=\sqrt{9^2+12^2}=\sqrt{225}=\res{15\text{ km}}$.

\ex{2} $\sqrt{6^2+0{,}8^2}=\sqrt{36{,}64}\approx\res{6{,}05\text{ m}}$.

\ex{3} $x^2+(x+2)^2=100\Rightarrow 2x^2+4x-96=0\Rightarrow x^2+2x-48=0\Rightarrow x=\frac{-2\pm14}{2}$: $x=6$ ($x=-8$ no tiene sentido). Catetos: $\res{6\text{ cm y }8\text{ cm}}$.

\ex{3} Si se rompe a altura $x$, el trozo caído mide $9-x$ y es la hipotenusa: $x^2+3^2=(9-x)^2\Rightarrow x^2+9=81-18x+x^2\Rightarrow18x=72\Rightarrow\res{x=4\text{ m}}$.


% =====================================================
\seccio{2}{Clasificación de triángulos por los lados}

\ex{1} \begin{tasks}(3)
  \task $169=169$: \rest{rectángulo}
  \task $64<85$: \rest{acutángulo}
  \task $81>52$: \rest{obtusángulo}
  \task $676=676$: \rest{rectángulo}
  \task $144>113$: \rest{obtusángulo}
  \task $121<181$: \rest{acutángulo}
\end{tasks}

\ex{1} \begin{tasks}(5)
  \task \rest{sí}
  \task no
  \task \rest{sí}
  \task \rest{sí}
  \task no
\end{tasks}($144+256=400$; $81+144=225\neq256$; $225+1296=1521$; $121+3600=3721$; $64+81=145\neq144$.)

\ex{2} \begin{tasks}(3)
  \task \rest{no}: $9>3+5$
  \task \rest{sí}, obtusángulo ($49>32$)
  \task \rest{no}: $12=5+7$
\end{tasks}

\ex{2} $2{,}25+4=6{,}25=2{,}5^2$: \rest{sí}, es rectángulo. \ $A=\frac{1{,}5\cdot2}{2}=\res{1{,}5\text{ m}^2}$.

\ex{2} $\sqrt{90^2+120^2}=\res{150\text{ cm}}$. Con 152 cm: $152^2>150^2$, el ángulo opuesto a la diagonal es \rest{obtuso} (más de 90°).

\ex{3} Si $x$ es la hipotenusa: $x=\sqrt{81+144}=\res{15}$. \ Si la hipotenusa es 12: $x=\sqrt{144-81}=\sqrt{63}\approx\res{7{,}94}$.

\ex{3} Comparamos $x^2$ con $36+64=100$, y además tiene que ser $x<6+8=14$. \ Acutángulo: $\res{x=9}$; \ rectángulo: $\res{x=10}$; \ obtusángulo: $\res{x=11,\,12,\,13}$.


% =====================================================
\seccio{3}{Aplicaciones del teorema de Pitágoras}

\ex{1} \begin{tasks}(3)
  \task $\sqrt{625}=\res{25\text{ cm}}$
  \task $\sqrt{98}\approx\res{9{,}90\text{ cm}}$
  \task $\sqrt{6{,}76}=\res{2{,}6\text{ m}}$
\end{tasks}

\ex{1} Isósceles: $h=\sqrt{17^2-8^2}=\sqrt{225}=\res{15\text{ cm}}$. \ Equilátero: $h=\sqrt{12^2-6^2}=\sqrt{108}\approx\res{10{,}39\text{ cm}}$.

\ex{1} $l=\sqrt{9^2+12^2}=\sqrt{225}=\res{15\text{ cm}}$; \ $P=4\cdot15=\res{60\text{ cm}}$.

\ex{2} $l^2+l^2=144\Rightarrow l=\sqrt{72}\approx\res{8{,}49\text{ cm}}$; \ $P=4\sqrt{72}\approx\res{33{,}94\text{ cm}}$.

\ex{2} $c=\sqrt{(15-9)^2+8^2}=\sqrt{100}=\res{10\text{ cm}}$; \ $P=15+9+8+10=\res{42\text{ cm}}$.

\ex{2} $\frac{22-12}{2}=5$; \ $h=\sqrt{13^2-5^2}=\sqrt{144}=\res{12\text{ cm}}$.

\ex{2} \begin{tasks}(3)
  \task $\sqrt{9+16}=\res{5}$
  \task $\sqrt{36+64}=\res{10}$
  \task $\sqrt{16+25}=\sqrt{41}\approx\res{6{,}40}$
\end{tasks}

\ex{3} $AB=6$, \ $AC=\sqrt{9+16}=5$, \ $BC=\sqrt{9+16}=5$: \rest{isósceles}, $P=\res{16}$. \ $36<25+25$: \rest{acutángulo}.

\ex{3} Mitad de la otra diagonal: $\sqrt{10^2-6^2}=8$; \ diagonal: $\res{16\text{ cm}}$.

\ex{3} $d=32\cdot2{,}54=81{,}28$ cm. \ $(16k)^2+(9k)^2=81{,}28^2\Rightarrow337k^2=6606{,}44\Rightarrow k\approx4{,}43$. \ Anchura $\approx\res{70{,}8\text{ cm}}$, altura $\approx\res{39{,}8\text{ cm}}$.


% =====================================================
\seccio{4}{Perímetros y áreas de triángulos y cuadriláteros}

\ex{1} \begin{tasks}(2)
  \task $A=\res{81\text{ cm}^2}$, $P=\res{36\text{ cm}}$
  \task $A=\res{60\text{ cm}^2}$, $P=\res{34\text{ cm}}$
  \task $A=\res{60\text{ cm}^2}$, $P=\res{34\text{ cm}}$
  \task $A=\res{96\text{ cm}^2}$; $l=\sqrt{64+36}=10$, $P=\res{40\text{ cm}}$
\end{tasks}

\ex{1} \begin{tasks}(2)
  \task $\frac{(11+7)\cdot5}{2}=\res{45\text{ m}^2}$
  \task $\frac{14\cdot9}{2}=\res{63\text{ cm}^2}$
\end{tasks}

\ex{1} \begin{tasks}(4)
  \task $\res{35\,000\text{ cm}^2}$
  \task $\res{4{,}5\text{ m}^2}$
  \task $\res{23\,000\text{ m}^2}$
  \task $\res{7{,}5\text{ m}^2}$
\end{tasks}

\ex{1} $h=\frac{72}{9}=\res{8\text{ cm}}$; \ $P=2\cdot9+2\cdot8=\res{34\text{ cm}}$.

\ex{2} $h=\sqrt{64-16}=\sqrt{48}\approx6{,}93$ cm; \ $A=\frac{8\cdot\sqrt{48}}{2}\approx\res{27{,}71\text{ cm}^2}$.

\ex{2} $h=\sqrt{100-36}=8$; \ $A=\frac{12\cdot8}{2}=\res{48\text{ cm}^2}$; \ $P=\res{32\text{ cm}}$.

\ex{2} $h=\sqrt{13^2-5^2}=12$; \ $A=\frac{(20+10)\cdot12}{2}=\res{180\text{ cm}^2}$; \ $P=20+10+2\cdot13=\res{56\text{ cm}}$.

\ex{2} $l=10$; mitad de la otra diagonal $\sqrt{100-64}=6$, \ $d=12$; \ $A=\frac{16\cdot12}{2}=\res{96\text{ cm}^2}$.

\ex{2} $A=\res{10\,200\text{ m}^2}=\res{1{,}02\text{ ha}}$; \ coste: $\res{30\,600\text{ €}}$.

\ex{3} $A=4{,}5\cdot2{,}6-1{,}2\cdot1=11{,}7-1{,}2=10{,}5\text{ m}^2$; \ $\frac{10{,}5}{4}=2{,}625$: \ $\res{3\text{ botes}}$.

\ex{3} $A=\frac{20\cdot15}{2}=\res{150\text{ cm}^2}$; \ $c=\sqrt{625}=\res{25\text{ cm}}$. Tomando la hipotenusa como base: $\frac{25\cdot h}{2}=150\Rightarrow\res{h=12\text{ cm}}$.

\ex{3} $b+h=17$ y $b^2+h^2=169$: \ $b^2+(17-b)^2=169\Rightarrow b^2-17b+60=0\Rightarrow b=12$ o $b=5$. Lados $\res{12\text{ cm y }5\text{ cm}}$; \ $A=\res{60\text{ cm}^2}$.


% =====================================================
\seccio{5}{Polígonos regulares}

\ex{1} \begin{tasks}(1)
  \task $P=\res{30\text{ cm}}$, \ $A=\frac{30\cdot4{,}13}{2}=\res{61{,}95\text{ cm}^2}$
  \task $P=\res{30\text{ cm}}$, \ $A=\frac{30\cdot4{,}33}{2}=\res{64{,}95\text{ cm}^2}$
  \task $P=\res{32\text{ cm}}$, \ $A=\frac{32\cdot4{,}83}{2}=\res{77{,}28\text{ cm}^2}$
\end{tasks}

\ex{1} $r=l=\res{8\text{ cm}}$; \ $a=\sqrt{64-16}=\sqrt{48}\approx\res{6{,}93\text{ cm}}$; \ $A=\frac{48\cdot\sqrt{48}}{2}\approx\res{166{,}28\text{ cm}^2}$.

\ex{2} $a=\sqrt{10{,}2^2-6^2}=\sqrt{68{,}04}\approx\res{8{,}25\text{ cm}}$; \ $A\approx\frac{60\cdot8{,}25}{2}\approx\res{247{,}5\text{ cm}^2}$.

\ex{2} La diagonal es un diámetro: $l^2+l^2=12^2\Rightarrow l=\sqrt{72}\approx\res{8{,}49\text{ cm}}$; \ $a=\frac l2\approx\res{4{,}24\text{ cm}}$; \ $A=l^2=\res{72\text{ cm}^2}$.

\ex{2} $a=\sqrt{100-25}=\sqrt{75}$; \ $A=\frac{60\cdot\sqrt{75}}{2}\approx259{,}81\text{ cm}^2$. \ $\frac{30\,000}{259{,}81}\approx115{,}5$: \ $\res{116\text{ baldosas}}$.

\ex{3} $a=2$; \ $\frac l2=\sqrt{4^2-2^2}=\sqrt{12}\approx3{,}46\Rightarrow l\approx\res{6{,}93\text{ cm}}$; \ $A=\frac{3\cdot2\sqrt{12}\cdot2}{2}\approx\res{20{,}78\text{ cm}^2}$.


% =====================================================
\seccio{6}{La circunferencia: longitud y arcos}

\ex{1} \begin{tasks}(3)
  \task $14\pi\approx\res{43{,}98\text{ cm}}$
  \task $15\pi\approx\res{47{,}12\text{ cm}}$
  \task $5\pi\approx\res{15{,}71\text{ m}}$
\end{tasks}

\ex{1} $r=\frac{62{,}83}{2\pi}\approx\res{10\text{ cm}}$; \ $d\approx\res{20\text{ cm}}$.

\ex{1} \begin{tasks}(3)
  \task $4\pi\approx\res{12{,}57\text{ cm}}$
  \task $2\pi\approx\res{6{,}28\text{ cm}}$
  \task $5\pi\approx\res{15{,}71\text{ cm}}$
\end{tasks}

\ex{2} Cada vuelta: $60\pi\approx188{,}50$ cm. \ $\frac{150\,000}{60\pi}\approx\res{795{,}8\text{ vueltas}}$ (unas 796).

\ex{2} Las dos semicircunferencias forman una circunferencia: $200+2\pi\cdot36{,}5=200+73\pi\approx\res{429{,}34\text{ m}}$.

\ex{2} 20 min son $\frac13$ de vuelta (120°): \ $\frac{2\pi\cdot15}{3}=10\pi\approx\res{31{,}42\text{ cm}}$.

\ex{3} $\frac{2\pi\cdot12\cdot n}{360}=18{,}85\Rightarrow n=\frac{18{,}85\cdot360}{24\pi}\approx\res{90^\circ}$.

\ex{3} $L=2\pi\cdot150\,000\,000\approx942\,477\,796$ km; \ 1 año $=365\cdot24=8760$ h; \ $v=\frac{L}{8760}\approx\res{107\,589\text{ km/h}}$.


% =====================================================
\seccio{7}{Área del círculo y figuras circulares}

\ex{1} \begin{tasks}(3)
  \task $25\pi\approx\res{78{,}54\text{ cm}^2}$
  \task $36\pi\approx\res{113{,}10\text{ cm}^2}$
  \task $0{,}64\pi\approx\res{2{,}01\text{ m}^2}$
\end{tasks}

\ex{1} \begin{tasks}(3)
  \task $6\pi\approx\res{18{,}85\text{ cm}^2}$
  \task $\frac{125\pi}{3}\approx\res{130{,}90\text{ cm}^2}$
  \task $3\pi\approx\res{9{,}42\text{ cm}^2}$
\end{tasks}

\ex{1} $\pi(81-16)=65\pi\approx\res{204{,}20\text{ cm}^2}$.

\ex{2} $r=\sqrt{\frac{200}{\pi}}\approx\res{7{,}98\text{ cm}}$; \ $L=2\pi r\approx\res{50{,}13\text{ cm}}$.

\ex{2} $\frac{\pi(144-49)\cdot80}{360}=\frac{190\pi}{9}\approx\res{66{,}32\text{ cm}^2}$.

\ex{2} $A=\frac{225\pi}{8}\approx\res{88{,}36\text{ cm}^2}$; \ arco: $\frac{30\pi}{8}\approx\res{11{,}78\text{ cm}}$.

\ex{2} $l=r=5$: \ circunscrito $25\pi\approx\res{78{,}54\text{ cm}^2}$. \ $a^2=25-6{,}25=18{,}75$: \ inscrito $18{,}75\pi\approx\res{58{,}90\text{ cm}^2}$.

\ex{3} Sector: $\frac{36\pi\cdot60}{360}=6\pi$. \ Triángulo equilátero de lado 6: $h=\sqrt{27}$, $A=3\sqrt{27}\approx15{,}59$. \ Segmento: $6\pi-3\sqrt{27}\approx\res{3{,}26\text{ cm}^2}$.

\ex{3} $\pi(100-r^2)=50\pi\Rightarrow r^2=50\Rightarrow r=\sqrt{50}\approx\res{7{,}07\text{ cm}}$.


% =====================================================
\seccio{8}{Figuras compuestas y problemas}

\ex{1} $A=6\cdot6+\frac{6\cdot4}{2}=36+12=\res{48\text{ m}^2}$. \ Lados del tejado: $\sqrt{3^2+4^2}=5$; \ $P=6+6+6+5+5=\res{28\text{ m}}$.

\ex{1} $A=64+\frac{\pi\cdot4^2}{2}=64+8\pi\approx\res{89{,}13\text{ cm}^2}$; \ $P=3\cdot8+\frac{2\pi\cdot4}{2}=24+4\pi\approx\res{36{,}57\text{ cm}}$.

\ex{2} Los dos semicírculos forman un círculo de radio 3: \ $A=60-9\pi\approx\res{31{,}73\text{ cm}^2}$; \ $P=2\cdot10+2\pi\cdot3=20+6\pi\approx\res{38{,}85\text{ cm}}$.

\ex{2} $A=144-\frac{\pi\cdot144}{4}=144-36\pi\approx\res{30{,}90\text{ cm}^2}$.

\ex{2} $A=1{,}4\cdot1{,}8+\frac{\pi\cdot0{,}7^2}{2}=2{,}52+0{,}245\pi\approx\res{3{,}29\text{ m}^2}$; \ $P=1{,}4+2\cdot1{,}8+0{,}7\pi\approx\res{7{,}20\text{ m}}$.

\ex{2} Rectángulo exterior: $(25+4)\times(12+4)=29\cdot16=464$; \ $A=464-300=\res{164\text{ m}^2}$.

\ex{3} Media hoja = cuarto de círculo $-$ medio cuadrado $=25\pi-50$. \ $A=2(25\pi-50)=50\pi-100\approx\res{57{,}08\text{ cm}^2}$.

\ex{3} Radio 5: \ $A=300-3\cdot25\pi=300-75\pi\approx\res{64{,}38\text{ cm}^2}$.

\ex{3} $r=4$, \ $a=\sqrt{16-4}=\sqrt{12}$; \ hexágono: $\frac{24\sqrt{12}}{2}=12\sqrt{12}\approx41{,}57$; \ $A=16\pi-12\sqrt{12}\approx50{,}27-41{,}57=\res{8{,}70\text{ cm}^2}$.

\ex{3} Tres cuartos de círculo de radio 5 y, al doblar las esquinas, dos cuartos de círculo de radio $5-4=1$: \ $A=\frac34\cdot25\pi+2\cdot\frac{\pi}{4}=19{,}25\pi\approx\res{60{,}48\text{ m}^2}$.


\end{document}
